\section*{v2.10.0 Trace Reconstruction and Observable Extraction Core}
\addcontentsline{toc}{section}{v2.10.0 Trace Reconstruction and Observable Extraction Core}

\begin{quote}
\textbf{Normalization note (v2.17.1 local review).}
This block is retained as a \emph{historical reconstruction and observable-extraction precursor}. Its reconstruction language remains useful, but the normalized proof-bearing architecture now lies in the unfolding phases and the grouped transport logic.
\end{quote}

The filtered-observability layer from v2.9.0 separates admissible modes into
readable, tagged, and mute sectors. The next structural task is no longer only to
screen admissible modes, but to reconstruct theorem-level observables from the
readable part without double counting, without reopening forbidden channels, and
without losing the ledger discipline enforced by HC confinement. The present step
therefore adds a trace-reconstruction layer. It answers the question: once a mode
has passed admissibility and readability tests, how is it promoted to an actual
observable in the companion proof architecture?

\subsection*{1. Reconstruction operator for readable modes}
Let $\{\mathcal F_\alpha\}_{\alpha\in A_{\mathrm{read}}}$ be the approved readable
filter family from v2.9.0. We now introduce a reconstruction operator
\[
\mathcal R: \bigoplus_{\alpha\in A_{\mathrm{read}}} \mathfrak X_\alpha \to \mathfrak O,
\]
where $\mathfrak X_\alpha = \mathcal F_\alpha \mathfrak X$ denotes the readable mode
sector and $\mathfrak O$ is the structural observable space. The role of
$\mathcal R$ is not arbitrary synthesis: it is constrained reconstruction. For a
state $u\in\mathfrak X$ one seeks a decomposition
\[
\mathrm{Obs}(u) = \mathcal R\bigl((\mathcal F_\alpha u)_{\alpha\in A_{\mathrm{read}}}\bigr)
\]
such that every observable component remains HC-confined, ledger-bounded, and
source-tagged whenever the reconstruction is not exact. Thus observable extraction
is treated as a disciplined readout map rather than an unrestricted projection.

\paragraph{Principle 4 (reconstructive observability).}
A readable mode contributes to theorem-level content only after passing through a
reconstruction operator whose output is compatible with HC confinement, channel
budgets, and source tags. Readability alone is therefore necessary but not yet
sufficient for observable status.

\subsection*{2. Observable extraction at theorem level}
For each admissible theorem-context $\Theta$ let $\mathfrak O_\Theta$ denote the class
of observables allowed in that context. One then requires extraction estimates of
schematic form
\[
\|\mathcal R_\Theta((\mathcal F_\alpha u)_\alpha)\|_{\mathfrak O_\Theta}
\le C_\Theta
\Bigl(
\sum_{\alpha\in A_{\mathrm{read}}} \mathfrak M_\alpha(u)
+ \mathrm{Tag}(u)+\mathrm{Leak}(u)
\Bigr),
\]
with $\mathfrak M_\alpha$ the readable trace masses from v2.9.0. This means an
observable is only legitimate when its theorem-level norm is controlled by the same
ledger quantities already used to regulate channel transport and spectral readout.
The companion proof therefore acquires an explicit extraction stage between readable
modes and theorem statements.

\subsection*{3. Overlap ledger and anti-double-counting rule}
Once multiple filters contribute to one observable quantity, overlap must be made
visible. Introduce an overlap ledger
\[
\mathrm{OL}(u)= \sum_{\alpha\neq\beta}
\omega_{\alpha\beta}\,\langle \mathcal F_\alpha u,\mathcal F_\beta u\rangle_{\mathfrak X},
\]
where the coefficients $\omega_{\alpha\beta}$ encode admissible correlation or
subtraction rules. The purpose of $\mathrm{OL}(u)$ is not to forbid overlap but to
prevent silent duplication of structural mass. Observable reconstruction is required
to satisfy a compensated estimate of the form
\[
\|\mathrm{Obs}(u)\|_{\mathfrak O}
\le C_{\mathrm{obs}}
\Bigl(
\sum_{\alpha\in A_{\mathrm{read}}} \mathfrak M_\alpha(u)
- \mathrm{OL}(u)
+ \mathrm{Tag}(u)+\mathrm{Leak}(u)
\Bigr),
\]
whenever the overlap ledger is well defined. Thus one can read several filtered
channels at once without pretending they were disjoint when they are not.

\subsection*{4. HC-stable observable closure}
The reconstructed observable sector must itself remain HC-stable. Concretely, one
requires that
\[
\Pi_{\mathrm{HC}}\,\mathrm{Obs}(u)=\mathrm{Obs}(u)+\tau(u),
\qquad
\|\tau(u)\|_{\mathfrak O}\le C_{\mathrm{HC}}
\bigl(\mathrm{Tag}(u)+\mathrm{Leak}(u)\bigr),
\]
so that reconstruction does not reactivate forbidden directions except through
ledger-visible correction terms. In this sense HC confinement becomes not only a
state-space condition and not only a filter-compatibility condition, but now also an
observable-closure condition. Only outputs that survive this test may be used as
stable structural readings.

\subsection*{5. Probe-sector reconstruction inheritance}
For the two-probe subsystem the same reconstruction principle must descend to the
probe readout. Let $\mathcal R^{\mathrm{probe}}$ be the probe reconstruction operator.
Then every probe observable should either be inherited from an ambient observable
through the descent chain or be explicitly tagged as probe-local correction. One
asks schematically for bounds of the form
\[
\|\mathrm{Obs}^{\mathrm{probe}}(v)\|_{\mathfrak O_{\mathrm{probe}}}
\le C_{\mathrm{pr}}
\bigl(
\|\mathrm{Obs}(u)\|_{\mathfrak O}
+\mathrm{Mismatch}(u,v)
\bigr),
\]
for matched ambient/probe states $(u,v)$. This ensures the probe sector remains an
instrument of structural descent rather than an independent source of untracked
observables.

\subsection*{6. Consequence for field equations and appendix tools}
The admission checklist for future candidate field equations is therefore extended
once more:
\begin{enumerate}[label=(\arabic*)]
  \item verify semigroup dissipativity and ledger closure;
  \item verify spectral admissibility and resolvent compatibility;
  \item verify readable mode separation under HC-compatible filters;
  \item construct a reconstruction operator for theorem-level observables; and
  \item bound overlap, leakage, and probe descent in a visible ledger form.
\end{enumerate}
The appendix tool taxonomy is refined accordingly. Besides resolvent estimators,
readout filters, and gap witnesses, a strong tool may now also be labelled as
reconstruction operator, observable extractor, overlap-ledger bound, HC-stability
witness, or probe-descent certificate. These labels mark tools that do not merely
prove admissibility, but actually promote admissible readable content into stable
observables.

\subsection*{7. v\docversion\ readiness criterion}
The present step should be regarded as successful if the following can now be done:
\begin{itemize}
  \item readable spectral content can be reconstructed into theorem-level observables;
  \item overlaps between readable channels are accounted for instead of silently duplicated;
  \item observable extraction remains HC-stable modulo visible tagged/leakage correction;
  \item probe observables descend from ambient observables rather than appearing autonomously; and
  \item future field equations can be screened for observable reconstruction, not merely for readable spectra.
\end{itemize}

This is the gain of v\docversion. The companion monolith now carries not only
closure, channel specialization, transport dynamics, forward invariance,
field-energy realization, semigroup readout, spectral admissibility, and filtered
observability, but also a first disciplined reconstruction layer. Future field
equations are therefore constrained four times: they must generate the correct
ledger-dissipative semigroup, belong to the correct resolvent class, pass a
trace-compatible readout filter, and finally admit HC-stable observable
reconstruction before they can count as theorem-level structural content.


